% \chapter*{あとがき}
\chapter*{Afterword}
% \section*{2版にむけたあとがき} 
\section*{Afterword for the Second Edition}

% 授業資料の印刷の手間を省くよ，という目的のためにKDPをやってみましたが，授業履修者以外の人にも手に取ってもらえたようで何よりです。まずは御礼申し上げます。
I tried using KDP to save the trouble of printing class materials, and I'm glad to see that people other than the students taking the class have picked it up. First of all, let me express my thanks.

% お礼に続いて謝罪になります。
Following my thanks, I would like to apologize.
% もう\textbf{ごめんなさい}としか言いようがないのですが，誤字脱字，誤変換が大量に含まれていたのみならず，動かないコードや「そうはならんやろ」という数式展開の間違い，引用した方のお名前のミスなどたくさんのたくさんの過ちをばらまいております。横山光輝の3国志だったかとおもいますが，諸葛孔明が「ふむぅ，さすが書物に嘘は書いてないものだ」という台詞がありまして，それを思い出しては「ごめん，俺のアレには嘘が書いてある・・・」と思って凹んでいます。授業中に気づいて「あれ？なんかこれおかしいな」と思うところは口頭で修正しているのですが，記録が残ってしまうというのは怖いですね。反省するだけじゃなくて，再発しないようにちゃんとしろってことなんですが，差し迫る授業期日にむけて必死で書いてるときは気づかないもので。
I can only say, "I'm sorry", but not only do I have numerous typos and mistranslations, but I also have code that doesn't work, errors in mathematical formula expansions that go against reason, and misspellings of the names of the people I've quoted. I've littered many such mistakes. I think it was Mitsuteru Yokoyama's Romance of the Three Kingdoms where Zhuge Liang says, "Hmm, as expected, there are no lies written in books", and I'm reminded of that line and think, "I'm sorry, there are lies in my documents..." and then I get downhearted. I try to correct what I notice is strange during my classes verbally, but the fact that it still gets recorded is scary. Reflecting is not enough, I need to make sure not to repeat these errors, but I don't notice them when I'm frantically writing for an impending class deadline.
% 皆様のご指摘を受けてサイト公開版ではバージョン1.2.20まで上がっています。常に最新版をチェックしていただけるとありがたく存じます。今回，かなり色々な問題点を修正したので，KDPの方でも版を改めることにしました。
Upon receiving feedback from all of you, the website's public version has been upgraded to version 1.2.20. I would appreciate it if you could always check the latest version. This time, since we've fixed quite a few issues, we've decided to revise the version on KDP as well.

% さて今回の大きな変更点としては，Stanのバージョンアップに合わせて\verb|array|表記のコードに変えたこと，\texttt{cmdstanr}パッケージをメイン環境に変えたことがあります。\texttt{rstan}パッケージは開発が遅く，実行の中断などもできないことがあって少し不便ですから，これからの(Rでの)Stanユーザは\texttt{cmdstanr}パッケージを使うことをおすすめします。問題は導入が少し面倒になったところでしょうか。パッケージをインストールした後に本体をインストールしたり，パスを設定したりする必要があります。またあの憎き最大手のOSが，初心者にわからないような仕掛けを仕組んでいるせいで，インストールするドライブがOneDriveでユーザ名にマルチバイト文字が入っているとエラーが出る，という問題が頻発しています。一応，対応を書いてくれているサイトに言及していますが，この手の情報はネットで検索してもらったほうが早く確実かもしれません。
The major changes this time are that we have changed the code to \verb|array| notation to match the upgrade of Stan, and have switched to the \texttt{cmdstanr} package as the main environment. The \texttt{rstan} package is slow to develop, and there are times when it cannot be interrupted, which is a little inconvenient, so we recommend that future (in R) Stan users use the \texttt{cmdstanr} package. The problem might be that the introduction has become a little more troublesome. After installing the package, you need to install the main body and set the path. There is also a major OS that has a mechanism that beginners do not understand, and if the drive to install is OneDrive and the username contains multibyte characters, an error will occur, which is a common problem. I refer to the site where the response is written, but it might be faster and more reliable to search for this information on the Internet.

% こうした問題はありますが，初学者は何も悪くないし，Stanでわかる世界は相変わらず楽しいです。みなさんフリーで楽しい心理統計モデリングの世界を楽しんでいっていただければと思います。
Though such problems exist, beginners aren't doing anything wrong and the world as understood by Stan continues to be enjoyable. I hope everyone can have fun in this free world of psychological statistical modeling.
\begin{flushright}
	2022-10-07\\
	小杉考司
\end{flushright}


% \section*{初版にむけたあとがき}
\section*{Afterword for the First Edition}

% このテキストは専修大学人間科学部心理学専攻2年次配当の講義，心理学データ解析応用のテキストです。1年次に姉妹講義であるデータ解析基礎を履修済みであることを想定していますが，単位未習得であってもわかるように心がけて書いているつもりです。
This text is for the lecture on Applied Psychological Data Analysis, assigned in the second year of the Department of Psychology, Faculty of Human Sciences, Senshu University. It is assumed that you have completed the sister lecture, Basic Data Analysis, in your first year, but even if you have not acquired the credits, it is intended to be written in a way that is understandable.

% 2020年から始まったコロナの問題で，2020年度の前期はオンライン，後期は一部対面での授業となりました。翌2021年度は逆に前期は対面，後期はオンラインが主の講義スタイルになりました。それまでは口頭でできていたことが，より明示的に静的な資料として事前に提示されている必要があると考え，授業内容の書き起こしをしていくことになりました。
Due to the coronavirus issue that began in 2020, the first half of the 2020 academic year was online, and the second half was a mix of face-to-face classes. In contrast, the 2021 academic year started with face-to-face classes in the first semester, while the second semester primarily consisted of online lectures. Because we had to give explicit instructions that could previously be done verbally, we found it necessary to present more static materials in advance. As a result, we started to transcribe the contents of the classes.
% データ解析基礎と同様，詳細なシラバスを準備してそれに対応するような授業教材を作ることにしました。基礎編と併せてご利用いただければと思います。
As with the basics of data analysis, we decided to prepare a detailed syllabus and create teaching materials to correspond to it. We hope you find it useful along with the basics.
% なお，付録など一部基礎編と内容が重複するところがあります。
Please note, there may be parts, such as the appendices, where the content overlaps with the basic content.

% 前半は多変量解析を，後半はベイジアンモデリングという2つのテーマがあるため，どちらかだけ履修することもあるでしょう。とはいえ，行列表現やRの使い方など，両者を通じての説明もありますからテキストとしては一冊にまとめました。
The course is divided into two themes: multivariate analysis in the first half, and Bayesian modeling in the second half. Therefore, you may choose to only take one of them. However, things like matrix representations and how to use R programming are explained in both sections, thus was compiled into one textbook.



% このテキストにはバージョン番号が付与されています。今回はバージョン1.2.1となっています。バージョン番号は，メジャー，マイナー，パッチの3つの要素から構成されています。
This text has a version number attached to it. This time it is version 1.2.1. The version number is composed of three elements: major, minor, and patch.
% 今はメジャーバージョン2です。
It's now major version 2.
% シラバスの骨組みが変われば，この数字が1つ増えることになります。
If the framework of the syllabus changes, this number will increase by one.
% チャプターレベルの変更があればマイナーの数字が上がります。
If there are changes at the chapter level, the minor number will increase.
% 誤字脱字など，細かな変更はパッチの数字が変わります。
Typos and minor changes will change the number of patches.
% なにせ自作のテキストですから，誤字脱字や誤変換，数字や参照のミスが多く，受講生に指摘してもらってはアップデートする，ということを繰り返してここまでやってきました。
After all, this is my original text, so there are many mistakes such as typos, mistranslations, incorrect numbers and references. I've come this far by repeatedly updating the text as students point out the mistakes.



% 最新版のシラバスやテキストは，サイト\footnote{\url{https://kosugitti.github.io/psychometrics_syllabus/}}で公開されています。
The latest syllabus and text are available on the website\footnote{\url{https://kosugitti.github.io/psychometrics_syllabus/}}.
% GitHubを使って管理し，図版のほとんどはRを使って作っています。版組にはもちろんLua\LaTeX を使っています。
We manage using GitHub, and most of the illustrations are created using R. Of course, we use Lua\LaTeX for typesetting.
% バージョンがこまめに上がるものですし，リンク機能などPDFならではの利点もあるので，できれば電子デバイスで電子ファイルのまま利用してもらえればと思っています。
The version is frequently upgraded, and there are advantages unique to PDFs, such as link functions, so I hope that you can use the electronic files on electronic devices if possible.
% しかし受講生の多くはプリントアウトして使っているようですから，この度実費で製本したものが入手できるよう，Kindle Direct Publishingを利用してペーパーバック版も用意しました。
However, since many students seem to be using it by printing out, this time we have prepared a paperback version using Kindle Direct Publishing so that you can obtain one at actual cost.
% 専修大学人間科学部心理学専攻の学生向けに用意したテキストではありますが，どなたでもお読みいただけるようになっています。内容についてはできるだけの注意を払って，誤解や間違いがないよう気をつけてはおりますが，おかしなところがありましたら著者までご連絡ください。
This text is prepared for the students of the Department of Psychology, Faculty of Human Sciences, Senshu University, but it is intended for anyone to read. While I have taken as much care as possible regarding the content and have strived to avoid misunderstandings or errors, if you find anything amiss, please contact the author.

% また原稿そのものはCreative Commons BY-SA(CC BY-SA)ライセンスVersion 4.0に基づいて提供しています。加筆修正や再頒布など，ライセンスにしたがってご活用いただければと思います。
The manuscript itself is provided under the Creative Commons BY-SA (CC BY-SA) License Version 4.0. We hope you will utilize it according to the license, such as by making additions, revisions, and redistributions.


\begin{flushright}
	2022-03-19\\
	小杉考司
\end{flushright}
